\chapter{The Step Function, the Principal Value, and the Missing \texorpdfstring{$\delta$}{delta}}
\label{ch:step}

\begin{quote}
\itshape
The derivative rules give an equation for the spectrum of the step, and the equation has infinitely
many solutions. The loss of information is not an accident of the method: multiplication by
$\omega$ annihilates $\delta$, exactly as differentiation in time annihilates constants. Recovering
the missing term requires a new object, the Cauchy principal value, and an independent computation.
\end{quote}

\section{The derivative of a step}

Let
\begin{equation}
H(t)=\begin{cases}0,&t<0,\\ 1,& t>0,\end{cases}
\label{eq:H}
\end{equation}
with the value at $0$ left unspecified. The choice at a single point will be irrelevant, and
Chapter~\ref{ch:dist} explains why.

\begin{proposition}\label{prop:Hprime}
As a distribution, $H'=\delta$.
\end{proposition}

\begin{proof}
For $\varphi\in C_{c}^{\infty}$,
\begin{equation}
\ip{H'}{\varphi}=-\ip{H}{\varphi'}=-\int_{0}^{\infty}\varphi'(t)\dd t
=-\big[\varphi(t)\big]_{0}^{\infty}=-\big(0-\varphi(0)\big)=\varphi(0)=\ip{\delta}{\varphi}.
\end{equation}
The boundary at $+\infty$ vanishes because $\varphi$ is compactly supported; the boundary at $0$
produces $\varphi(0)$.
\end{proof}

The mechanism is worth stating separately, since it is the source of every $\delta$ in applications.

\begin{proposition}[piecewise smooth functions]\label{prop:jump}
Let $f$ be $C^{1}$ except at points $a_{j}$ accumulating only at infinity, with jumps
$J_{j}=f(a_{j}^{+})-f(a_{j}^{-})$ and a classical derivative $f'_{\mathrm{cl}}$ off the jump set.
Then
\begin{equation}
Df=f'_{\mathrm{cl}}+\sum_{j}J_{j}\,\delta_{a_{j}},
\label{eq:jump}
\end{equation}
where $\delta_{a}(\varphi)=\varphi(a)$.
\end{proposition}

\begin{proof}
Split the integral $\int f\varphi'$ over the intervals between the jumps and integrate by parts on
each. The interior terms reassemble into $-\int f'_{\mathrm{cl}}\varphi$, and each jump $a_{j}$
contributes $\big(f(a_{j}^{+})-f(a_{j}^{-})\big)\varphi(a_{j})=J_{j}\ip{\delta_{a_{j}}}{\varphi}$.
\end{proof}

\begin{example}\label{ex:abs}
$D|t|=\sgn t$ and $D^{2}|t|=2\delta$. The first has a jump of size $2$ in the derivative of
$|t|$ at the origin; the second differentiates that jump.
\end{example}

\section{Applying the differentiation rule}

By \eqref{eq:derivt}, applied to $H$ in the distributional sense and to $H'=\delta$,
\begin{equation}
i\omega\,\hf(\omega)=\F[H'](\omega)=\F[\delta](\omega)=1 .
\label{eq:iomegaH}
\end{equation}
Read off naively this suggests $\hf(\omega)=1/(i\omega)$. There are two independent reasons why
that cannot be the end of the story.

\begin{enumerate}
\item $1/\omega$ is not locally integrable at $0$: $\int_{0}^{1}\omega^{-1}\dd\omega=\infty$. It
therefore does not define a distribution by integration, and ``divide by $\omega$'' is not an
operation available in this setting.
\item Even granting a meaning to $1/\omega$ near the origin, the equation $i\omega\,\hf=1$ does not
determine $\hf$: if $U$ is one solution then $U+c\delta$ is another, because multiplication by
$\omega$ annihilates $\delta$.
\end{enumerate}

\begin{proposition}\label{prop:omega-delta}
$\omega\delta=0$, that is, $\ip{\omega\delta}{\varphi}=\ip{\delta}{\omega\varphi}=0$ for every
test function $\varphi$.
\end{proposition}

\begin{proof}
By definition of multiplication by a smooth function,
$\ip{\omega\delta}{\varphi}=\ip{\delta}{\omega\varphi}=0\cdot\varphi(0)=0$.
\end{proof}

\begin{remark}[the same loss in both domains]
In the time domain, differentiating kills constants: $(H+C)'=H'$ for every $C$. In the frequency
domain, multiplying by $\omega$ kills $\delta$: $\omega(U+c\delta)=\omega U$ for every $c$. These are
one phenomenon viewed twice. The differentiation rule can therefore only determine the part of the
spectrum away from $\omega=0$; whatever mass sits exactly at zero frequency is invisible to it.
\end{remark}

\section{The principal value}

The function $1/\omega$ has a non-integrable singularity at the origin, but symmetric truncation
removes it. That observation defines a distribution.

\begin{definition}[Cauchy principal value]\label{def:PV}
For $\varphi\in C_{c}^{\infty}(\R)$ set
\begin{equation}
\left\langle\PV\frac{1}{\omega},\varphi\right\rangle
:=\lim_{a\downarrow0}\int_{|\omega|>a}\frac{\varphi(\omega)}{\omega}\dd\omega .
\label{eq:PVdef}
\end{equation}
\end{definition}

The limit can be computed in closed form by pairing the two half lines.

\begin{lemma}\label{lem:PVform}
For every $\varphi\in C_{c}^{\infty}(\R)$ and every $a>0$,
\begin{equation}
\int_{|\omega|>a}\frac{\varphi(\omega)}{\omega}\dd\omega
=\int_{a}^{\infty}\frac{\varphi(x)-\varphi(-x)}{x}\dd x .
\label{eq:PVpaired}
\end{equation}
Consequently
\begin{equation}
\left\langle\PV\frac{1}{\omega},\varphi\right\rangle
=\int_{0}^{\infty}\frac{\varphi(x)-\varphi(-x)}{x}\dd x .
\label{eq:PVint}
\end{equation}
\end{lemma}

\begin{proof}
In the integral over $\omega<-a$ substitute $\omega=-x$ and add the two pieces.
\end{proof}

The point of \eqref{eq:PVpaired} is the numerator: $\varphi(x)-\varphi(-x)$ vanishes at $x=0$ at
first order, so dividing by $x$ leaves a bounded integrand. Each half line separately would
diverge; the pair does not.

\begin{theorem}[existence]\label{thm:PVexists}
The functional \eqref{eq:PVint} is well defined on $C_{c}^{\infty}(\R)$ and is a distribution.
\end{theorem}

\begin{proof}
Let $\varphi\in C_{c}^{\infty}$ and let $K$ be a compact interval containing its support. Split
$\int_{0}^{\infty}$ at $x=1$.

For $0<x\le1$, the mean value theorem gives
\begin{equation}
|\varphi(x)-\varphi(-x)|\le 2x\sup_{|y|\le1}|\varphi'(y)| ,
\qquad\text{hence}\qquad
\left|\frac{\varphi(x)-\varphi(-x)}{x}\right|\le 2\|\varphi'\|_{\infty,[-1,1]} ,
\end{equation}
which is integrable on $(0,1]$.

For $x\ge1$, compact support gives $\varphi(\pm x)=0$ once $x$ exceeds the support radius, so the
integrand vanishes there; in any case $|\varphi(x)-\varphi(-x)|\le 2\|\varphi\|_{\infty}$ and
$1/x$ is integrable on $[1,\infty)$ intersected with the support after bounding.

Linearity is clear. Continuity follows from the displayed estimate together with the bound on
$[1,\infty)$, which together give
\begin{equation}
\left|\left\langle\PV\tfrac{1}{\omega},\varphi\right\rangle\right|
\le C_{K}\Big(\|\varphi'\|_{\infty,[-1,1]}+\|\varphi\|_{\infty}\Big). \qedhere
\end{equation}
\end{proof}

\begin{remark}[the symmetry is part of the definition]\label{rem:PVasym}
It is not true that any way of removing the singularity gives the same answer. Consider the
asymmetric truncation that keeps $a<|\omega|<1$ on the right and $b<|\omega|<1$ on the left, with
$a,b\downarrow0$. By \eqref{eq:PVpaired} and a split at $b$,
\begin{equation}
\int_{a<|\omega|<1}\frac{\varphi(\omega)}{\omega}\dd\omega
=\int_{a}^{1}\frac{\varphi(x)-\varphi(-x)}{x}\dd x+\int_{a}^{b}\frac{\varphi(-x)}{x}\dd x ,
\end{equation}
and the second term equals $\varphi(0)\log(b/a)+o(1)$. Hence the limit exists only when
$b/a\to1$, and then equals the principal value; otherwise the answer differs by
$\varphi(0)\log(b/a)$. The principal value is a specific symmetric limit, not the assertion that
the singularity does not matter.
\end{remark}

Two elementary properties are needed later.

\begin{proposition}\label{prop:PVprops}
$\PV\frac{1}{\omega}$ is an odd distribution, and
\begin{equation}
\omega\,\PV\frac{1}{\omega}=1 .
\label{eq:omegaPV}
\end{equation}
\end{proposition}

\begin{proof}
Oddness follows from \eqref{eq:PVint}: replacing $\varphi$ by its reflection reverses the sign of the
integrand. For \eqref{eq:omegaPV}, use the definition of multiplication by $\omega$ and
\eqref{eq:PVint} with $\omega\varphi$ in place of $\varphi$:
\begin{equation}
\left\langle\omega\PV\tfrac{1}{\omega},\varphi\right\rangle
=\left\langle\PV\tfrac{1}{\omega},\omega\varphi\right\rangle
=\int_{0}^{\infty}\frac{x\varphi(x)-(-x)\varphi(-x)}{x}\dd x
=\int_{0}^{\infty}\big(\varphi(x)+\varphi(-x)\big)\dd x=\int_{\R}\varphi . \qedhere
\end{equation}
\end{proof}

\section{The regularized step, computed}

The natural regularization of $H$ is one-sided: decay only where $H$ is nonzero. Set
\begin{equation}
H_{\ep}(t)=H(t)e^{-\ep t},\qquad \ep>0 .
\label{eq:Heps}
\end{equation}

\begin{theorem}\label{thm:Heps}
For every $\ep>0$ and $\omega\in\R$,
\begin{equation}
\F[H_{\ep}](\omega)=\frac{1}{\ep+i\omega}
=\frac{\ep}{\ep^{2}+\omega^{2}}-i\,\frac{\omega}{\ep^{2}+\omega^{2}} .
\label{eq:Hepssplit}
\end{equation}
\end{theorem}

\begin{proof}
Directly,
\begin{equation}
\int_{0}^{\infty}e^{-(\ep+i\omega)t}\dd t
=\left[\frac{e^{-(\ep+i\omega)t}}{-(\ep+i\omega)}\right]_{0}^{\infty}
=\frac{1}{\ep+i\omega},
\end{equation}
the upper boundary vanishing because $\ep>0$; then rationalize the denominator.
\end{proof}

\begin{remark}[why one-sided]
With the two-sided regularization $e^{-\ep|t|}$ the limit is $1$
(Chapter~\ref{ch:constant}); with the one-sided one it is $H$. The shape of the decay determines the
limit object, not merely the speed of convergence. This is the first place where the choice of
regularization carries information.
\end{remark}

The two terms of \eqref{eq:Hepssplit} have to be treated separately: the real part is even and
concentrates, the imaginary part is odd and must be paired across the origin.

\section{The two limits}

\begin{theorem}[real part]\label{thm:realpart}
\begin{equation}
\frac{\ep}{\ep^{2}+\omega^{2}}\longrightarrow\pi\delta(\omega)
\qquad(\ep\downarrow0),
\label{eq:realpart}
\end{equation}
the limit being taken in the sense of test functions.
\end{theorem}

\begin{proof}
Since $\dfrac{\ep}{\ep^{2}+\omega^{2}}=\pi K_{\ep}$ with $K_{\ep}$ as in \eqref{eq:Keps}, and since
$\ip{K_{\ep}}{\varphi}\to\varphi(0)$ for every bounded continuous $\varphi$
(Theorem~\ref{thm:sampling}), one gets
$\ip{\pi K_{\ep}}{\varphi}\to\pi\varphi(0)=\ip{\pi\delta}{\varphi}$.
\end{proof}

The coefficient $\pi$ is the total mass:
\begin{equation}
\int_{\R}\frac{\ep}{\ep^{2}+\omega^{2}}\dd\omega=\pi ,
\label{eq:masspi}
\end{equation}
independently of $\ep$. Whenever a family of even functions concentrates, its limit is the mass
times $\delta$; this is the single most useful heuristic in the subject, and here it is
exactly what the theorem says.

\begin{theorem}[imaginary part]\label{thm:imagpart}
\begin{equation}
\frac{\omega}{\ep^{2}+\omega^{2}}\longrightarrow\PV\frac{1}{\omega}
\qquad(\ep\downarrow0).
\label{eq:imagpart}
\end{equation}
\end{theorem}

\begin{proof}
Let $\varphi\in C_{c}^{\infty}$. Pair the half lines by substituting $\omega=-x$ in the negative
one:
\begin{equation}
\int_{\R}\frac{\omega}{\ep^{2}+\omega^{2}}\varphi(\omega)\dd\omega
=\int_{0}^{\infty}\frac{x}{\ep^{2}+x^{2}}\big[\varphi(x)-\varphi(-x)\big]\dd x .
\label{eq:imagsplit}
\end{equation}
We apply dominated convergence to \eqref{eq:imagsplit} and must exhibit a dominating function
independent of $\ep$.

\emph{Near the origin}, for $0<x\le1$, use $x/(\ep^{2}+x^{2})\le1/x$ together with the mean value
estimate $|\varphi(x)-\varphi(-x)|\le2x\sup_{|y|\le1}|\varphi'|$:
\begin{equation}
\left|\frac{x}{\ep^{2}+x^{2}}\big[\varphi(x)-\varphi(-x)\big]\right|
\le 2\sup_{|y|\le1}|\varphi'(y)| ,
\end{equation}
a constant, hence integrable on $(0,1]$.

\emph{Far from the origin}, for $x\ge1$, the same inequality $x/(\ep^{2}+x^{2})\le1/x$ combined
with compact support gives a bound $C/x$ on a bounded interval, hence integrable on $[1,\infty)$.

For each fixed $x>0$, $x/(\ep^{2}+x^{2})\to1/x$ as $\ep\downarrow0$. Dominated convergence in
\eqref{eq:imagsplit} therefore yields
\begin{equation}
\int_{\R}\frac{\omega}{\ep^{2}+\omega^{2}}\varphi(\omega)\dd\omega
\longrightarrow\int_{0}^{\infty}\frac{\varphi(x)-\varphi(-x)}{x}\dd x
=\left\langle\PV\frac{1}{\omega},\varphi\right\rangle ,
\end{equation}
which is \eqref{eq:imagpart}.
\end{proof}

\begin{remark}[why the two regions must be separated]
The kernel $\omega/(\ep^{2}+\omega^{2})$ is small near the origin and small at infinity, with a
maximum of $1/(2\ep)$ at $\omega=\pm\ep$. A single global estimate cannot capture both ends; the
split at $x=1$ makes the near part bounded by the oscillation of $\varphi$ and the far part bounded
by the decay of $\varphi$. This ``control near the singular point by smoothness, control at
infinity by decay'' is the standard pattern, and it recurs verbatim in
Chapter~\ref{ch:limits}.
\end{remark}

\section{The spectrum of the step}

Combining Theorems~\ref{thm:Heps}, \ref{thm:realpart} and \ref{thm:imagpart}, and using
$1/i=-i$,
\begin{equation}
\F[H](\omega)=\pi\delta(\omega)-i\,\PV\frac{1}{\omega}
=\pi\delta(\omega)+\PV\frac{1}{i\omega} .
\label{eq:FH}
\end{equation}

The coefficient $\pi$ in front of $\delta$ deserves an independent verification, since it was
obtained from a computation rather than from a structural argument.

\begin{theorem}[symmetry determines the coefficient]\label{thm:symcoef}
Let $c\in\C$ and suppose $\F[H]=c\,\delta-i\,\PV\frac{1}{\omega}$. Then $c=\pi$.
\end{theorem}

\begin{proof}
As distributions, $H(t)+H(-t)=1$: the identity holds for every $t\neq0$, and the value at a single
point does not affect the distribution. Applying $\F$ and using
$\F[f(-\cdot)](\omega)=\fh(-\omega)$,
\begin{equation}
\hf(\omega)+\hf(-\omega)=2\pi\delta(\omega).
\label{eq:Hsym}
\end{equation}
Substituting the ansatz and using that $\PV\frac{1}{\omega}$ is odd while $\delta$ is even,
\begin{equation}
\big(c\delta-i\PV\tfrac{1}{\omega}\big)+\big(c\delta+i\PV\tfrac{1}{\omega}\big)=2c\,\delta .
\end{equation}
Comparing with \eqref{eq:Hsym} gives $2c=2\pi$, hence $c=\pi$.
\end{proof}

\begin{corollary}[the sign function]\label{cor:sgn}
\begin{equation}
\F[\sgn t](\omega)=-2i\,\PV\frac{1}{\omega}=\frac{2}{i}\,\PV\frac{1}{\omega} .
\label{eq:sgn}
\end{equation}
\end{corollary}

\begin{proof}
Since $H=\tfrac12+\tfrac12\sgn$ and $\F[1]=2\pi\delta$,
\begin{equation}
\F[\sgn t]=2\Big(\F[H]-\tfrac12\F[1]\Big)
=2\big(\pi\delta-i\PV\tfrac{1}{\omega}-\pi\delta\big)
=-2i\,\PV\tfrac{1}{\omega}. \qedhere
\end{equation}
\end{proof}

\begin{remark}[why no $\delta'$ can be added]
The equation $i\omega\hf=1$ is compatible with adding $c\delta$, since $\omega\delta=0$, but not
with adding $\delta'$, since $\omega\delta'=-\delta\neq0$
(\eqref{eq:xdeltan}). Proposition~\ref{prop:keromega} makes this precise: the solutions of
$\omega T=1$ differ exactly by multiples of $\delta$. So \eqref{eq:FH} is the general solution, and
the only question was the value of $c$, settled by \eqref{eq:Hsym}.
\end{remark}

\section{Boundary values}

The computation \eqref{eq:Hepssplit} also produces the standard notation of mathematical physics.
For $\ep\downarrow0$,
\begin{equation}
\frac{1}{x\pm i0}:=\lim_{\ep\downarrow0}\frac{1}{x\pm i\ep}
=\PV\frac{1}{x}\mp i\pi\,\delta(x),
\label{eq:boundary}
\end{equation}
where the limit is distributional. The notation is a shorthand for the limit, not a substitution of
a complex zero.

To see the sign in \eqref{eq:FH} by this route, note that
\begin{equation}
\frac{1}{\ep+i\omega}=\frac{-i}{\omega-i\ep}\longrightarrow\frac{-i}{\omega-i0}
=-i\Big(\PV\frac{1}{\omega}+i\pi\delta\Big)=\pi\delta-i\,\PV\frac{1}{\omega},
\end{equation}
which reproduces \eqref{eq:FH}. The two minus signs---one from $1/i=-i$, one from the lower sign in
\eqref{eq:boundary}---cancel in the $\delta$ term and add in the $\PV$ term; this is the most
common place for a sign error to enter, and it is worth doing the computation twice.

\section{The transform of the principal value}

One more table entry can be obtained for free from what has already been proved, and it illustrates how
the reflection identity \eqref{eq:F2} is used in practice.

\begin{proposition}\label{prop:PVtransform}
In $\Sw'$,
\begin{equation}
\F\Big[\PV\frac{1}{t}\Big](\omega)=-i\pi\,\sgn\omega .
\label{eq:PVtransform}
\end{equation}
\end{proposition}

\begin{proof}
Apply $\F^{2}T=2\pi\,T(-\cdot)$ to $T=\sgn t$. By \eqref{eq:sgn},
$\F[\sgn t]=-2i\PV(1/\omega)$, so
\begin{equation}
\F\Big[\F[\sgn t]\Big]=\F\Big[-2i\,\PV\frac{1}{\omega}\Big]=-2i\,\F\Big[\PV\frac{1}{\omega}\Big]
\end{equation}
while the reflection identity gives
$\F^{2}[\sgn t]=2\pi\,\sgn(-t)=-2\pi\,\sgn t$. Equating,
\begin{equation}
-2i\,\F\Big[\PV\frac{1}{\omega}\Big]=-2\pi\,\sgn\omega ,
\qquad\text{hence}\qquad
\F\Big[\PV\frac{1}{\omega}\Big]=\frac{\pi}{i}\,\sgn\omega=-i\pi\,\sgn\omega . \qedhere
\end{equation}
\end{proof}

\begin{remark}[a consistency check on the factor]
The answer must be odd and purely imaginary, because $\PV(1/t)$ is odd and real; it is. A sharper
check is to apply $\F$ once more and use \eqref{eq:F2}. On one hand
\begin{equation}
\F\Big[\F\Big[\PV\frac{1}{t}\Big]\Big]=\F\big[-i\pi\,\sgn\omega\big]=-i\pi\cdot\big(-2i\big)\PV\frac{1}{\omega}
=-2\pi\,\PV\frac{1}{\omega},
\end{equation}
using \eqref{eq:sgn}; on the other hand \eqref{eq:F2} gives
$2\pi\,\PV\frac{1}{-t}=-2\pi\,\PV\frac{1}{t}$. The two agree, as they must, and the factor
$\pi$ is consistent with \eqref{eq:Fcos}.
\end{remark}

\section{Summary of the chapter}

\begin{enumerate}
\item $H'=\delta$, and more generally each jump of size $J$ contributes $J\delta$ to the
distributional derivative (Propositions~\ref{prop:Hprime}, \ref{prop:jump}).
\item The differentiation rule yields $i\omega\hf=1$, which does not determine $\hf$ because
$\omega\delta=0$ (Proposition~\ref{prop:omega-delta}).
\item $1/\omega$ is not locally integrable; its symmetric regularization is the principal value,
which is a well-defined distribution (Theorems~\ref{thm:PVexists}) with
$\omega\PV\frac{1}{\omega}=1$.
\item The asymmetric truncation shows that the symmetry in the definition is essential
(Remark~\ref{rem:PVasym}).
\item Regularizing the step one-sidedly and separating real and imaginary parts gives
$\pi\delta$ and $-i\PV\frac{1}{\omega}$ (Theorems~\ref{thm:realpart}, \ref{thm:imagpart}).
\item The coefficient $\pi$ is forced by $H(t)+H(-t)=1$ (Theorem~\ref{thm:symcoef}).
\item Consequently $\F[H]=\pi\delta+\PV\frac{1}{i\omega}$ and
$\F[\sgn t]=-2i\PV\frac{1}{\omega}$.
\end{enumerate}
